%==============================
\section{Case Study}
\label{sec:ch04_case_study}

The top-down coordination framework is illustrated next on a
representative multi-robot mission with global temporal structure,
hierarchical team organization, and online task updates. The purpose of
the case study is to show how decomposition, assignment, team
formation, and replanning interact in a single execution flow, and to
provide empirical context for the complexity discussion of the
preceding sections.

Consider a mission specified initially at the team level through a global
temporal formula that combines ordered service requirements, allocation
constraints, and limited concurrency. In the first stage, the accepted behavior
of the task automaton is pruned and converted into subtask occurrences together
with a compact precedence structure. This relaxed partially ordered view
separates the logical dependencies that must be preserved from the robot-level
details that can be optimized later. The resulting structure is then used to
assign subtasks to heterogeneous robots while minimizing the mission-level
completion metric. At this stage, the role of top-down synthesis is explicit:
the global specification is not approximated by ad hoc task splitting, but is
systematically refined into allocable components whose temporal relations remain
traceable to the original formula.

In the second stage, the assigned subtasks are grouped into teams and refined
through the hierarchical fleet-coordination machinery of
Section~\ref{ch04:subsec:task}. High-level decomposition determines which robots
should cooperate, team formation enforces capability and redundancy conditions,
and local coordination within each team resolves the motion-level realization of
the assigned subtask sequence. The hierarchy becomes especially valuable when
the fleet is large or heterogeneous, because it prevents the global mission from
being reintroduced as a single fully coupled planning object at every level.
Instead, each layer exposes only the information required by the next one.

Finally, suppose that a new mission element or operator request arrives during
execution. The current global plan is then updated through the online product of
R-posets, and the receding-horizon mechanism of
Section~\ref{ch04:sec:online_continual_tasks} revises only the unexecuted suffix
of the plan. Tasks already in progress are preserved whenever possible, while
new obligations are incorporated into the evolving precedence structure. This
case-study view makes clear that the top-down framework is not merely an
offline decomposition procedure. It is a persistent coordination architecture in
which global temporal structure, hierarchical allocation, and online revision
remain coupled throughout execution.

%==============================
\section{Summary}
\label{ch04:sec:summary}

This chapter developed a top-down framework for global temporal-task
coordination in multi-robot systems. Section~\ref{ch04:sec:po_decomp_assignment}
introduced automaton pruning, task decomposition, relaxed partially ordered
sets, and branch-and-bound assignment for a fixed global mission.
Section~\ref{ch04:subsec:task} extended this framework to large heterogeneous
fleets through hierarchical mission decomposition, abstract team assignment,
team formation, and local coordination within teams.
Section~\ref{ch04:sec:online_continual_tasks} then addressed continual tasks by
updating the task structure online through R-poset products and by applying a
receding-horizon replanning scheme. Human requests were incorporated as special
online task updates within the same framework.
